% GENERATED FILE. Edit canonical lessons, then run npm run generate:books.
% canonical-source-hash: fnv1a64:0ff71cc5fecd5224
% canonical-lessons: RU-C430-postel, RU-C430-prostynya, RU-C430-navolochka, RU-C430-gornichnaya, RU-C430-stelit

\chapter{Making the bed}
\label{ch:ru-making-the-bed}
\hlchaptermodality{430}

\begin{chapteropening}
\textbf{By the end of this chapter:} \emph{I can talk about a hotel bed: bedding, a sheet, a pillowcase, the chambermaid and making up a bed.}

Complete the last lesson of chapter 430: I can talk about a hotel bed: bedding, a sheet, a pillowcase, the chambermaid and making up a bed.
\end{chapteropening}

\section[\texorpdfstring{\ru{постель} (postél)}{постель (postél)}]{\ru{постель} (postél) — bedding; a bed}

\label{lesson:RU-C430-postel}



\begin{quote}
Before the new one: say the Russian for a deposit, then the Russian for to extend.
\end{quote}

\subsection*{You'll want to know: \ru{постель}}
\textbf{\ru{постель}} (\emph{postél}) — \textquotedblleft{}bedding; a bed\textquotedblright{}.

\begin{quote}
\ru{Постель} \ru{в} \ru{номере} \ru{была} \ru{чистая} \ru{и} \ru{мягкая}.
\end{quote}

\begin{quote}
\emph{Postél v nómere bylá chístaya i myágkaya.}
\end{quote}

\begin{quote}
The bed in the room was clean and soft.
\end{quote}

Feminine. \emph{\ru{Постельное} \ru{бельё}} — bed linen.

\begin{cousinweb}
\emph{\ru{постель}} is built on the root of \emph{\ru{стлать}}, \emph{\ru{стелить}} 'to spread': what is spread out to sleep on. EtymDB lists \emph{\ru{постель}} with the senses 'bed; bedding'.
\end{cousinweb}

\begin{grammarlens}[title={What you've built}]
One more word for telling, explaining and describing at B1.
\end{grammarlens}

\subsection*{Your turn}
\begin{itemize}
\raggedright
  \item \emph{Say it:} \emph{postél}
  \item \emph{Say it:} \emph{postél}, once more
  \item \emph{Say it:} a sentence of your own with \emph{postél}
  \item \emph{Recall:} say the Russian for a deposit, then the Russian for to extend, then say \emph{postél} again
\end{itemize}

\begin{tcolorbox}[breakable,colback=teal!4,colframe=teal!35!black,title={Before you move on}]
What does \ru{постель} mean? (\textquotedblleft{}Bedding; a bed\textquotedblright{}.) Say it once more.
\end{tcolorbox}
\section[\texorpdfstring{\ru{простыня} (prostynyá)}{простыня (prostynyá)}]{\ru{простыня} (prostynyá) — a sheet (on a bed)}

\label{lesson:RU-C430-prostynya}



\begin{quote}
Before the new one: say the Russian for to extend, then the Russian for bedding.
\end{quote}

\subsection*{You'll want to know: \ru{простыня}}
\textbf{\ru{простыня}} (\emph{prostynyá}) — \textquotedblleft{}a sheet (on a bed)\textquotedblright{}.

\begin{quote}
\ru{В} \ru{шкафу} \ru{лежат} \ru{чистые} \ru{простыни}.
\end{quote}

\begin{quote}
\emph{V shkafú lezhát chístyye próstyni.}
\end{quote}

\begin{quote}
There are clean sheets in the wardrobe.
\end{quote}

Feminine; plural \emph{\ru{простыни}}, genitive plural \emph{\ru{простынь}}.

\begin{cousinweb}
\emph{\ru{простыня}} is usually linked with \emph{\ru{простой}} 'plain, simple': a plain piece of cloth. EtymDB does not list the word, so the line is hedged.
\end{cousinweb}

\begin{grammarlens}[title={What you've built}]
One more word for telling, explaining and describing at B1.
\end{grammarlens}

\subsection*{Your turn}
\begin{itemize}
\raggedright
  \item \emph{Say it:} \emph{prostynyá}
  \item \emph{Say it:} \emph{prostynyá}, once more
  \item \emph{Say it:} a sentence of your own with \emph{prostynyá}
  \item \emph{Recall:} say the Russian for to extend, then the Russian for bedding, then say \emph{prostynyá} again
\end{itemize}

\begin{tcolorbox}[breakable,colback=teal!4,colframe=teal!35!black,title={Before you move on}]
What does \ru{простыня} mean? (\textquotedblleft{}A sheet (on a bed)\textquotedblright{}.) Say it once more.
\end{tcolorbox}
\section[\texorpdfstring{\ru{наволочка} (návolochka)}{наволочка (návolochka)}]{\ru{наволочка} (návolochka) — a pillowcase}

\label{lesson:RU-C430-navolochka}



\begin{quote}
Before the new one: say the Russian for bedding, then the Russian for a sheet.
\end{quote}

\subsection*{You'll want to know: \ru{наволочка}}
\textbf{\ru{наволочка}} (\emph{návolochka}) — \textquotedblleft{}a pillowcase\textquotedblright{}.

\begin{quote}
\ru{На} \ru{подушке} \ru{не} \ru{было} \ru{наволочки}.
\end{quote}

\begin{quote}
\emph{Na podúshke né bylo návolochki.}
\end{quote}

\begin{quote}
There was no pillowcase on the pillow.
\end{quote}

Feminine; the stress is on the first syllable.

\begin{cousinweb}
\emph{\ru{наволочка}} is built from \emph{\ru{на}-} 'onto' and the root of \emph{\ru{волочь}} 'to drag, to pull', which EtymDB takes from Proto-Slavic \emph{velťi} 'to draw': a cover pulled on.
\end{cousinweb}

\begin{grammarlens}[title={What you've built}]
One more word for telling, explaining and describing at B1.
\end{grammarlens}

\subsection*{Your turn}
\begin{itemize}
\raggedright
  \item \emph{Say it:} \emph{návolochka}
  \item \emph{Say it:} \emph{návolochka}, once more
  \item \emph{Say it:} a sentence of your own with \emph{návolochka}
  \item \emph{Recall:} say the Russian for bedding, then the Russian for a sheet, then say \emph{návolochka} again
\end{itemize}

\begin{tcolorbox}[breakable,colback=teal!4,colframe=teal!35!black,title={Before you move on}]
What does \ru{наволочка} mean? (\textquotedblleft{}A pillowcase\textquotedblright{}.) Say it once more.
\end{tcolorbox}
\section[\texorpdfstring{\ru{горничная} (górnichnaya)}{горничная (górnichnaya)}]{\ru{горничная} (górnichnaya) — a chambermaid, a maid}

\label{lesson:RU-C430-gornichnaya}



\begin{quote}
Before the new one: say the Russian for a sheet, then the Russian for a pillowcase.
\end{quote}

\subsection*{You'll want to know: \ru{горничная}}
\textbf{\ru{горничная}} (\emph{górnichnaya}) — \textquotedblleft{}a chambermaid, a maid\textquotedblright{}.

\begin{quote}
\ru{Каждое} \ru{утро} \ru{горничная} \ru{убирает} \ru{номер}.
\end{quote}

\begin{quote}
\emph{Kázhdoye útro górnichnaya ubiráyet nómer.}
\end{quote}

\begin{quote}
Every morning the chambermaid cleans the room.
\end{quote}

An adjective used as a noun; it declines like an adjective.

\begin{cousinweb}
\emph{\ru{горничная}} is built on \emph{\ru{горница}} 'a room', first an upper room. EtymDB traces \emph{\ru{горница}} to Old East Slavic \emph{gorĭnĭ}, a word that goes with \emph{\ru{гора}} 'a mountain'.
\end{cousinweb}

\begin{grammarlens}[title={What you've built}]
One more word for telling, explaining and describing at B1.
\end{grammarlens}

\subsection*{Your turn}
\begin{itemize}
\raggedright
  \item \emph{Say it:} \emph{górnichnaya}
  \item \emph{Say it:} \emph{górnichnaya}, once more
  \item \emph{Say it:} a sentence of your own with \emph{górnichnaya}
  \item \emph{Recall:} say the Russian for a sheet, then the Russian for a pillowcase, then say \emph{górnichnaya} again
\end{itemize}

\begin{tcolorbox}[breakable,colback=teal!4,colframe=teal!35!black,title={Before you move on}]
What does \ru{горничная} mean? (\textquotedblleft{}A chambermaid, a maid\textquotedblright{}.) Say it once more.
\end{tcolorbox}
\section[\texorpdfstring{\ru{стелить} (stelít)}{стелить (stelít)}]{\ru{стелить} (stelít) — to make up (a bed); to spread}

\label{lesson:RU-C430-stelit}



\begin{quote}
Before the new one: say the Russian for a pillowcase, then the Russian for a chambermaid.
\end{quote}

\subsection*{You'll want to know: \ru{стелить}}
\textbf{\ru{стелить}} (\emph{stelít}) — \textquotedblleft{}to make up (a bed); to spread\textquotedblright{}.

\begin{quote}
\ru{Бабушка} \ru{стелит} \ru{нам} \ru{постель} \ru{на} \ru{диване}.
\end{quote}

\begin{quote}
\emph{Bábushka stélit nam postél na diváne.}
\end{quote}

\begin{quote}
Grandma makes up a bed for us on the sofa.
\end{quote}

Imperfective; perfective \emph{\ru{постелить}}. Present \emph{\ru{стелю}, \ru{стелешь}}. Also \emph{\ru{стлать}}.

\begin{cousinweb}
\emph{\ru{стелить}} is a form of \emph{\ru{стлать}} 'to spread', the root of \emph{\ru{постель}} 'bedding'. EtymDB lists \emph{\ru{стелить}} without a deeper source.
\end{cousinweb}

\begin{grammarlens}[title={What you've built}]
One more word for telling, explaining and describing at B1.
\end{grammarlens}

\subsection*{Your turn}
\begin{itemize}
\raggedright
  \item \emph{Say it:} \emph{stelít}
  \item \emph{Say it:} \emph{stelít}, once more
  \item \emph{Say it:} a sentence of your own with \emph{stelít}
  \item \emph{Recall:} say the Russian for a pillowcase, then the Russian for a chambermaid, then say \emph{stelít} again
\end{itemize}

\begin{tcolorbox}[breakable,colback=teal!4,colframe=teal!35!black,title={Before you move on}]
What does \ru{стелить} mean? (\textquotedblleft{}To make up (a bed); to spread\textquotedblright{}.) Say it once more.
\end{tcolorbox}
